% ============================================================
% DOCUMENTCLASS
% ============================================================
\documentclass[12pt,oneside]{book}

% ============================================================
% METADATOS
% ============================================================
\newcommand{\universidad}{Universidad de Buenos Aires (UBA)}
\newcommand{\facultad}{Facultad de Derecho}
\newcommand{\materia}{Constitucional Profundizado}
\newcommand{\titulo}{Apuntes de Constitucional Profundizado}
\newcommand{\autor}{Isaac Edgar Camacho Ocampo}
\newcommand{\anio}{2026}

% ============================================================
% PAQUETES
% ============================================================
\usepackage[utf8]{inputenc}
\usepackage[T1]{fontenc}
\usepackage[spanish]{babel}
\usepackage[
    top=2.5cm,
    bottom=2.5cm,
    left=3cm,
    right=2.5cm,
    headheight=15pt
]{geometry}
\usepackage{amsmath}
\usepackage{amssymb}
\usepackage{mathtools}
\usepackage{graphicx}
\usepackage{float}
\usepackage{caption}
\usepackage{subcaption}
\usepackage{booktabs}
\usepackage{array}
\usepackage{longtable}
\usepackage{tabularx}
\usepackage{enumitem}
\usepackage{xcolor}
\usepackage[most]{tcolorbox}
\usepackage{fancyhdr}
\usepackage{ragged2e}
\usepackage[
    colorlinks=true,
    linkcolor=blue!50!black,
    citecolor=green!40!black,
    urlcolor=blue!60!black,
    pdftitle={\titulo},
    pdfauthor={\autor},
    pdfsubject={Apuntes universitarios},
    bookmarks=true
]{hyperref}

% ============================================================
% CONFIGURACIÓN
% ============================================================
\graphicspath{
    {./img/}
    {./graficos/}
    {./figuras/}
}
\setcounter{tocdepth}{2}
\setlength{\parindent}{0pt}
\setlength{\parskip}{0.5em}
\setlist[itemize]{
    topsep=0.3em,
    itemsep=0.2em
}
\setlist[enumerate]{
    topsep=0.3em,
    itemsep=0.2em
}
\clubpenalty=10000
\widowpenalty=10000

\pagestyle{fancy}
\fancyhf{}
\fancyhead[L]{\nouppercase{\leftmark}}
\fancyhead[R]{\materia}
\fancyfoot[C]{\thepage}
\renewcommand{\headrulewidth}{0.4pt}
\renewcommand{\footrulewidth}{0pt}
\fancypagestyle{plain}{
    \fancyhf{}
    \fancyfoot[C]{\thepage}
    \renewcommand{\headrulewidth}{0pt}
}

% ============================================================
% COMANDOS
% ============================================================
\newcommand{\abs}[1]{\left|#1\right|}
\newcommand{\norm}[1]{\left\lVert#1\right\rVert}
\newcommand{\vect}[1]{\mathbf{#1}}
\newcommand{\deriv}[2]{\frac{d#1}{d#2}}
\newcommand{\pderiv}[2]{\frac{\partial#1}{\partial#2}}

% Anchos de las columnas del cuadro Cornell
\newlength{\cornellizq}
\newlength{\cornellder}
\setlength{\cornellizq}{0.27\textwidth}
\setlength{\cornellder}{0.63\textwidth}

% ============================================================
% ENTORNOS / CAJAS
% ============================================================
\newtcolorbox{definition}[1]{
    enhanced,
    breakable,
    title={Definición: #1},
    fonttitle=\bfseries,
    colback=blue!3,
    colframe=blue!50!black,
    boxrule=0.8pt,
    arc=2mm
}

\newtcolorbox{important}{
    enhanced,
    breakable,
    title={Importante},
    fonttitle=\bfseries,
    colback=yellow!8,
    colframe=orange!70!black,
    boxrule=0.8pt,
    arc=2mm
}

\newtcolorbox{example}[1]{
    enhanced,
    breakable,
    title={Ejemplo: #1},
    fonttitle=\bfseries,
    colback=green!4,
    colframe=green!40!black,
    boxrule=0.8pt,
    arc=2mm
}

\newtcolorbox{formula}[1]{
    enhanced,
    breakable,
    title={Fórmula / Regla: #1},
    fonttitle=\bfseries,
    colback=purple!4,
    colframe=purple!50!black,
    boxrule=0.8pt,
    arc=2mm
}

\newtcolorbox{exercise}[1]{
    enhanced,
    breakable,
    title={Ejercicio: #1},
    fonttitle=\bfseries,
    colback=gray!5,
    colframe=gray!60!black,
    boxrule=0.8pt,
    arc=2mm
}

\newtcolorbox{note}{
    enhanced,
    breakable,
    title={Nota},
    fonttitle=\bfseries,
    colback=white,
    colframe=gray!50,
    boxrule=0.6pt,
    arc=2mm
}

% ============================================================
% CONFIGURACIÓN FINAL
% ============================================================
\renewcommand{\arraystretch}{1.2}

% ============================================================
% DOCUMENT
% ============================================================
\begin{document}

% ------------------------------------------------------------
% PORTADA
% ------------------------------------------------------------
\begin{titlepage}

\thispagestyle{empty}

\begin{center}

\vspace*{1cm}

{\large \textsc{\universidad}}

\vspace{0.2cm}

{\large \textsc{\facultad}}

\vspace{3cm}

{\Huge \textbf{\titulo}}

\vspace{0.6cm}

{\large Segunda parte de la cursada · Segundo cuatrimestre 2026}

\vspace{0.8cm}

\textit{Comprender los conceptos es el primer paso para convertir el estudio en conocimiento.}

\vspace{1.2cm}

\rule{0.70\textwidth}{0.5pt}

\vspace{0.5cm}

{\large Apuntes de estudio}

\vspace{0.5cm}

\rule{0.70\textwidth}{0.5pt}

\vfill

{\large \autor}

\vspace{0.3cm}

{\large \anio}

\end{center}

\vfill

\begin{flushleft}
\textit{Este es un modesto aporte a los estudiantes de ``Constitucional Profundizado'' no pretende sustituir el material oficial de la materia.}
\end{flushleft}

\end{titlepage}

% ------------------------------------------------------------
% ÍNDICE
% ------------------------------------------------------------
\tableofcontents

% ------------------------------------------------------------
% INTRODUCCIÓN
% ------------------------------------------------------------
\chapter*{Presentación}
\addcontentsline{toc}{chapter}{Presentación}
\markboth{Presentación}{Presentación}

Estos apuntes corresponden a una clase de repaso y afianzamiento de la segunda parte de la materia. Recorren cinco ejes: (1) la \textbf{defensa de la competencia}, que protege el mercado frente a los abusos de quien lo domina; (2) el \textbf{medio ambiente}, con sus principios y el reparto de competencias entre la Nación y las provincias; (3) la \textbf{propiedad y la expropiación}, es decir, el límite del Estado frente a los bienes de las personas; (4) la \textbf{responsabilidad del Estado} por su actividad legítima o ilegítima y su relación con los contratos; y (5) la \textbf{legitimación activa} según el caso Halabi, que ordena quién puede reclamar y mediante qué acción.

Los cinco ejes comparten una idea común: el poder tiene límites y la Constitución ofrece herramientas para controlarlos.

\begin{example}{La terminal automotriz que quiere instalarse en el país}
Los cinco bloques pueden vincularse mediante el caso de una empresa terminal automotriz que desea instalarse en el país:
\begin{itemize}
    \item \textbf{Competencia:} quiere comprar a un competidor local; la Autoridad de Defensa de la Competencia debe decidir si lo autoriza, lo condiciona o lo rechaza.
    \item \textbf{Ambiente:} necesita instalar una planta; debe cumplirse la evaluación de impacto ambiental, los seguros y la remediación.
    \item \textbf{Propiedad:} necesita un terreno y el Estado necesita una autopista; se plantea cuándo puede expropiar y cuánto debe pagar.
    \item \textbf{Responsabilidad estatal:} si el Estado cambia las reglas o falla en un servicio, se plantea quién responde y por cuánto.
    \item \textbf{Legitimación:} si la planta fabrica 25.000 camionetas con frenos defectuosos, se plantea quién puede reclamar y cómo.
\end{itemize}
\end{example}

% ============================================================
% CAPÍTULO 1
% ============================================================
\chapter{Defensa de la competencia}

\section{Reconocimiento constitucional y posición dominante}

La defensa de la competencia tiene hoy un reconocimiento y una protección constitucionales. Constituye un objetivo que debe primar en las políticas de gobierno, de modo de promover la libre competencia entre los distintos actores y sujetos económicos del mercado, y de prevenir las distorsiones que se producen generalmente a raíz de alguna posición dominante que pueda tener una parte en un mercado.

\begin{definition}{Posición dominante}
Situación en la cual una parte que brinda un bien o un servicio a consumidores dentro del mercado no está expuesta a una competencia sustancial, es decir, a una competencia real. Puede originarse por varios motivos.
\end{definition}

\begin{important}
La posición dominante, y en particular los monopolios, no se evalúan igual que la competencia real: quien no enfrenta competencia sustancial puede abusar de su situación, y por ello el Estado debe prevenir las distorsiones del mercado.
\end{important}

\section{Monopolios legales y naturales}

Uno de los orígenes posibles de la posición dominante es el monopolio, que puede ser legal o natural.

\subsection{El monopolio legal}

\begin{definition}{Monopolio legal}
Monopolio que se otorga a una empresa para operar con exclusividad un servicio dentro de una determinada zona. Se da generalmente en el ámbito de las concesiones de servicios públicos, es decir, de los servicios domiciliarios.
\end{definition}

Constituyen ejemplos de servicios prestados en exclusividad dentro de una zona geográfica:
\begin{itemize}
    \item el transporte y la distribución de energía (Edenor y Edesur);
    \item el servicio de agua y cloacas en el Gran Buenos Aires, que corresponde al ámbito de la antigua Obras Sanitarias de la Nación (según se expuso, al momento de la clase se encontraba en un proceso de reprivatización);
    \item el transporte y distribución de gas (por ejemplo, Metrogas);
    \item el servicio aeroportuario: en el ámbito de la Capital Federal hay un solo aeropuerto activo, Aeroparque, prestado por una empresa que tiene un virtual monopolio por ser prestador único.
\end{itemize}

Bajo la idea de servicio público se encierran, entonces, los monopolios legales: son servicios prestados por un prestador único, pero con una base legal que lo admite. En contraprestación por ser el único prestador y no estar expuesto a una competencia sustancial, la ley le impone una serie de obligaciones y prevé una fuerte intensidad regulatoria. Lo más importante es que se le \textbf{fija la tarifa}: al no existir una competencia al lado, no hay un poder disciplinador sobre la tarifa que se cobra por el servicio.

El marco legal que otorga la concesión prevé, entre otras cuestiones:
\begin{itemize}
    \item inversiones para el mantenimiento y la extensión de la red, en cualquiera de los servicios domiciliarios;
    \item disposiciones sobre la calidad de los servicios;
    \item la continuidad de los servicios, que en principio son continuos y no pueden ser interrumpidos;
    \item la universalidad del servicio;
    \item la fijación de la tarifa que retribuye el servicio.
\end{itemize}

Mediante la técnica jurídica del servicio público, que es la regulación más intensiva de un servicio, el Estado combate la posición dominante del prestador único, para que no se valga de ella y cobre lo que desee. Esto es así porque se trata de un monopolio legal, es decir, porque la propia ley ha admitido la prestación en exclusividad.

\subsection{El monopolio natural}

\begin{definition}{Monopolio natural}
Situación en la cual existe un recurso escaso y un único prestador que lo explota.
\end{definition}

% TODO: contenido incompleto o ambiguo en las notas originales (el régimen del monopolio natural solo se enuncia, no se desarrolla)
El régimen aplicable al monopolio natural solo fue enunciado en el apunte, sin mayor desarrollo.

\begin{table}[H]
\centering
\small
\caption{Monopolio legal y monopolio natural}
\begin{tabularx}{\textwidth}{>{\raggedright\arraybackslash\bfseries}p{2.6cm} >{\raggedright\arraybackslash}X >{\raggedright\arraybackslash}X}
\toprule
Aspecto & Monopolio legal & Monopolio natural \\
\midrule
Origen & La propia ley admite la prestación en exclusividad de un servicio dentro de una zona. & Existe un recurso escaso y un único prestador que lo explota. \\
Ámbito típico & Concesiones de servicios públicos: energía (Edenor, Edesur), gas (Metrogas), agua y cloacas, aeropuertos (Aeroparque). & Situaciones de recurso escaso con un solo explotador. \\
Cómo se combate la posición dominante & Regulación intensiva: se fijan inversiones, control de calidad, continuidad, universalidad y, en lo esencial, la tarifa. & Solo se enuncia; no se desarrolla su régimen. \\
\bottomrule
\end{tabularx}
\end{table}

\begin{important}
Quien es prestador único carece de un «poder disciplinador» de la competencia sobre su tarifa; por eso el Estado lo reemplaza con la técnica del servicio público: la regulación más intensiva de un servicio.
\end{important}

\section{Herramientas de la ley: preventiva y sancionatoria}

La ley procura evitar las distorsiones del mercado y los abusos de la posición dominante. Para ello establece dos grandes herramientas: una \textbf{preventiva} y una \textbf{sancionatoria}. Ambas están a cargo de la \textbf{Autoridad de Defensa de la Competencia}, órgano creado por la ley de defensa de la competencia (Ley 27.442), que tiene dos competencias:
\begin{itemize}
    \item la \textbf{fase sancionatoria}: fijar multas cuando se corrobora una infracción a la ley, por ejemplo el abuso de posición dominante;
    \item la \textbf{fase preventiva}: actuar de forma previa para evitar situaciones de excesiva concentración en los mercados.
\end{itemize}

\subsection{Fase preventiva: el control de las concentraciones}

Cuando una empresa es comprada, o cuando en un mercado un prestador compra a otro y puede llegar a configurarse un monopolio, interviene la Autoridad. Puede establecer que la compra entraña un peligro para la competencia por una excesiva concentración. Frente a ello puede tomar tres decisiones: \textbf{convalidar} la concentración económica, \textbf{condicionarla} o \textbf{rechazarla}.

\begin{table}[H]
\centering
\small
\caption{Decisiones de la Autoridad ante una concentración económica}
\begin{tabularx}{\textwidth}{>{\raggedright\arraybackslash\bfseries}p{4.2cm} >{\raggedright\arraybackslash}X}
\toprule
Decisión & Efecto concreto \\
\midrule
Convalidar (autorizar sin condicionamientos) & El comprador toma posesión y control de la nueva empresa y puede ampliarla. \\
Condicionar & El nuevo dueño debe cumplir los deberes que le impone la Autoridad de Defensa de la Competencia. \\
Rechazar & No puede hacerse el negocio jurídico. \\
\bottomrule
\end{tabularx}
\end{table}

Muchas adquisiciones, fusiones y transferencias de empresas se dan a nivel global. En cada país donde se desarrollan concentraciones económicas derivadas de grandes transacciones internacionales, la dependencia de defensa de la competencia local analiza la situación concreta de ese país.

\begin{example}{Compra de Quilmes por un grupo belga}
Un grupo belga muy grande, el más poderoso en el mercado, compró la cervecería argentina Quilmes. Ese grupo económico realizó compras a nivel global, por lo que en cada país la autoridad de competencia analiza la situación local.
\end{example}

\subsection{Fase sancionatoria: las prácticas anticompetitivas}

La fase sancionatoria opera cuando existen prácticas anticompetitivas. Los ejemplos expuestos son:
\begin{itemize}
    \item la \textbf{compra atada}, por ejemplo cuando en un supermercado se establece una promoción por la cual hay que comprar más productos para acceder a un beneficio;
    \item la \textbf{cartelización}: todos se ponen de acuerdo para fijar el precio de un bien;
    \item el acuerdo entre dos o más empresas para distribuirse zonas geográficas donde no hay competencia, de modo que cada una, dentro de su zona, pueda hacer lo que quiere. Se señaló que esta práctica es más común de lo que se cree.
\end{itemize}

\begin{table}[H]
\centering
\small
\caption{Las dos herramientas de la ley}
\begin{tabularx}{\textwidth}{>{\raggedright\arraybackslash\bfseries}p{2.6cm} >{\raggedright\arraybackslash}X >{\raggedright\arraybackslash}X}
\toprule
Aspecto & Fase preventiva & Fase sancionatoria \\
\midrule
Momento & Antes: se actúa previamente para evitar la excesiva concentración. & Después: se corrobora una infracción a la ley. \\
Qué hace la Autoridad & Analiza compras, fusiones y transferencias de empresas; aprueba, condiciona o rechaza. & Fija multas. \\
Ejemplos & Compra de una empresa por un competidor (grupo belga y Quilmes). & Compra atada en un supermercado; cartelización (acuerdo de precios); reparto de zonas geográficas; abuso de posición dominante. \\
\bottomrule
\end{tabularx}
\end{table}

\section{La Autoridad de Defensa de la Competencia}

\subsection{Conformación}

El órgano de la ley antitrust tiene \textbf{cinco miembros}. Al momento de la clase ya había tres designados. Se cumplió un hito importante: por primera vez se prestó el acuerdo del Senado a tres de sus miembros. Por ello, por primera vez desde hace más de 20 años, contados desde la ley anterior, el órgano llega a consolidarse como un tribunal independiente y autónomo, porque sus miembros tienen estabilidad en el cargo, otorgada por el acuerdo del Senado. Esto constituye un paso adelante en la institucionalización de esta oficina, que interviene en cualquier mercado y, por lo tanto, en toda la economía en general.

\begin{important}
La independencia del órgano depende de la estabilidad de sus miembros, que se logra con el acuerdo del Senado. La Autoridad interviene en cualquier mercado, es decir, en toda la economía.
\end{important}

\subsection{Designación y mandato}

Los miembros del tribunal:
\begin{itemize}
    \item son cinco;
    \item son designados por el Ejecutivo con el acuerdo del Congreso (acuerdo del Senado), previo concurso de antecedentes;
    \item duran cinco años en el cargo;
    \item se renuevan de manera escalonada, de modo que un mismo gobierno no termine de designar toda la autoridad administrativa: cada período de dos años y medio debe renovarse parte de los vocales.
\end{itemize}

En la ley anterior el órgano se denominaba Comisión de Defensa de la Competencia; hoy se denomina Autoridad de Defensa de la Competencia. La ley vigente es la 27.442.

\subsection{Las decisiones son apelables}

La Autoridad no tiene la última palabra sobre una sanción ni sobre una concentración. Si la parte lo desea, puede apelar, porque se trata de un organismo administrativo, que se desenvuelve dentro de la órbita de la administración pública; por lo tanto, sus decisiones no pueden ser definitivas. Las decisiones son apelables ante el fuero civil y comercial federal.

\begin{important}
La Autoridad es un organismo administrativo: no tiene la última palabra. Toda sanción o decisión sobre una concentración puede apelarse ante el fuero civil y comercial federal.
\end{important}

\section{Análisis del mercado y concurrencia}

Para decidir si existe una situación de posición dominante, la Autoridad debe identificar el mercado. Cada mercado es particular y se trata de cuestiones técnicas muy precisas.

\begin{example}{Bimbo y Fargo}
La multinacional mexicana Bimbo compró Fargo, que era el gran jugador del pan lactal. Se plantea cuál es el mercado relevante: si todos los panes, es decir, todas las panificaciones que consume la población, o una parte de ellas. En el caso se entendió que los panificados del estilo Fargo eran particulares: son alimentos menos perecederos, que tienden a durar más en el consumo para una determinada utilidad, pero existen sustitutos.
\end{example}

La Autoridad debe analizar, entonces, cuál es el mercado y si ese mercado tiene sustitutos. Debe analizar también las \textbf{barreras de ingreso y egreso}: cuanto más libre es el mercado y menos requisitos hay para ingresar o salir de él, más flexible puede ser el análisis y mayor el grado de concentración admisible. En cambio, cuando las barreras son altas hay que cuidarse mucho de la concentración.

\begin{example}{Medicamentos}
En el mercado de los medicamentos hay grandes barreras de ingreso: los medicamentos están sujetos a patentes y los laboratorios que los producen deben tener una determinada certificación y control. Por ello hay que cuidarse mucho de la concentración en el mercado. En mercados más dinámicos, con menos barreras de ingreso y egreso, se puede ser más flexible al analizar el grado de concentración.
\end{example}

\begin{table}[H]
\centering
\small
\caption{Qué analiza la Autoridad de Defensa de la Competencia}
\begin{tabularx}{\textwidth}{>{\raggedright\arraybackslash\bfseries}p{2.8cm} >{\raggedright\arraybackslash}X >{\raggedright\arraybackslash}X}
\toprule
Criterio & Qué se examina & Ejemplo \\
\midrule
Mercado relevante & Cada mercado es particular: hay que identificar cuál es. & Bimbo--Fargo: ¿todas las panificaciones o solo los panificados tipo pan lactal, menos perecederos y de mayor duración? \\
Sustitutos & Si el producto tiene sustitutos disponibles. & Para los panificados del estilo Fargo existen sustitutos. \\
Barreras de ingreso y egreso & A más barreras, más cuidado con la concentración; a menos barreras, mayor flexibilidad. & Medicamentos: patentes, certificación y control (barreras altas). Mercados dinámicos: barreras bajas. \\
\bottomrule
\end{tabularx}
\end{table}

\subsection{La concurrencia como objetivo del Estado}

Promover la libre competencia debería ser un objetivo del Estado en todas sus políticas. Se denomina \textbf{concurrencia} a la posibilidad de que dos o más empresas ofrezcan bienes en un mismo mercado. Esto se aplica también a la propia actividad del Estado: cuando el Estado compra o llama a una licitación, lo mejor es que se presente la mayor cantidad posible de empresas, porque así compiten entre sí y el Estado puede obtener el mejor precio por el bien o servicio. La materia es de amplio espectro y atraviesa toda la economía.

\section{Ejemplos de conductas anticompetitivas}

Son ejemplos de conducta antitrust o anticompetitiva:
\begin{itemize}
    \item concertar precios;
    \item restringir el crecimiento de otra empresa, por ejemplo mediante la regulación de cantidades de materias primas o de ventas;
    \item la negativa de venta;
    \item imponer condicionamientos que pongan en peligro la permanencia de otros rivales o competidores.
\end{itemize}

\begin{example}{Conducta predatoria de precios}
Una empresa se presenta en el mercado y, por un período muy prolongado, mantiene precios excesivamente bajos para que su competidor, que no tiene la misma espalda financiera, se funda. Así se queda con todo el mercado y luego sube los precios a lo que desea. Es una típica conducta predatoria de precios, castigada por la ley.
\end{example}

\section{Cuadro Cornell: defensa de la competencia}

\begin{longtable}{
    >{\raggedright\arraybackslash}p{\cornellizq}
    >{\raggedright\arraybackslash}p{\cornellder}
}
\toprule
\textbf{Preguntas / palabras clave} & \textbf{Notas / respuestas} \\
\midrule
\endhead

\textbf{¿Qué es una posición dominante y qué respaldo constitucional tiene la defensa de la competencia?} &
Es la situación en la que quien brinda un bien o servicio no está expuesto a una competencia sustancial o real. La defensa de la competencia tiene reconocimiento constitucional y es un objetivo que debe primar en las políticas de gobierno: promover la libre competencia y prevenir distorsiones. \\[0.8em]

\textbf{¿Qué diferencia hay entre monopolio legal y monopolio natural?} &
\textbf{Legal:} exclusividad que otorga la ley, típica de las concesiones de servicios públicos (energía, agua y cloacas, gas, aeropuertos); el Estado lo compensa con regulación intensiva: fija tarifa, inversiones, calidad, continuidad y universalidad. \textbf{Natural:} un recurso escaso explotado por un único prestador. \\[0.8em]

\textbf{¿Por qué se fija la tarifa del prestador único?} &
Porque, al no tener competencia, no existe un poder disciplinador sobre la tarifa; el Estado lo reemplaza con la técnica del servicio público. \\[0.8em]

\textbf{¿Cuáles son las dos herramientas de la ley y qué decisiones puede tomar la Autoridad?} &
\textbf{Preventiva:} controla concentraciones económicas y puede convalidarlas, condicionarlas o rechazarlas. \textbf{Sancionatoria:} aplica multas ante prácticas anticompetitivas (compra atada, cartelización, reparto de zonas, abuso de posición dominante). Las operaciones globales se analizan país por país. \\[0.8em]

\textbf{¿Cómo se conforma la Autoridad y qué la hace independiente?} &
Tiene cinco miembros, designados por el Ejecutivo con acuerdo del Senado y concurso de antecedentes, con cinco años de mandato y renovación escalonada. La estabilidad en el cargo, dada por el acuerdo del Senado, la consolida como órgano independiente y autónomo. \\[0.8em]

\textbf{¿Qué analiza la Autoridad para determinar si hay posición dominante?} &
El \textbf{mercado relevante} (caso Bimbo--Fargo), los \textbf{sustitutos} y las \textbf{barreras de ingreso y egreso} (medicamentos: barreras altas; mercados dinámicos: más flexibilidad). La concurrencia también aplica al Estado comprador: más oferentes en una licitación, mejor precio. \\[0.8em]

\textbf{¿Quién tiene la última palabra sobre una sanción y qué conductas se castigan?} &
Las decisiones son apelables ante el fuero civil y comercial federal, porque la Autoridad es un organismo administrativo. Se castigan, entre otras, concertar precios, restringir cantidades, negarse a vender, condicionar la permanencia de rivales y los precios predatorios. Ley 27.442. \\

\midrule
\multicolumn{2}{p{\dimexpr\cornellizq+\cornellder+2\tabcolsep\relax}}{%
\textbf{Resumen:} La defensa de la competencia, de reconocimiento constitucional, busca evitar las distorsiones del mercado y el abuso de la posición dominante. Frente a los monopolios legales, el Estado aplica regulación intensiva, en especial la fijación de tarifa. La Autoridad de Defensa de la Competencia previene (controla concentraciones: convalida, condiciona o rechaza) y sanciona (multas por prácticas anticompetitivas). Es un órgano de cinco miembros, independiente por la estabilidad que da el acuerdo del Senado, pero sus decisiones son apelables ante el fuero civil y comercial federal. } \\
\bottomrule
\end{longtable}

% ============================================================
% CAPÍTULO 2
% ============================================================
\chapter{Medio ambiente}

\section{Competencias ambientales: presupuestos mínimos}

Todas las personas tienen derecho a un ambiente sano y equilibrado. El daño ambiental es pasible de sanciones en muchos casos y genera el deber de reparación.

En materia de competencias, la Constitución prevé que regula el ambiente el Estado nacional, y que las provincias lo hacen de manera complementaria. Rigen conceptos de \textbf{complementariedad} y de \textbf{solidaridad} en la protección del ambiente.

\begin{definition}{Presupuestos mínimos ambientales}
Ley sancionada por el Congreso de la Nación que establece una base de obligaciones en materia ambiental: no se puede estar por debajo de esos presupuestos ni de esas obligaciones.
\end{definition}

Las provincias, por encima de esos presupuestos mínimos, pueden dictar normas en el ámbito de sus respectivas jurisdicciones para aumentar las previsiones y los requisitos de la protección ambiental.

\begin{table}[H]
\centering
\small
\caption{Reparto de competencias ambientales}
\begin{tabularx}{\textwidth}{>{\raggedright\arraybackslash\bfseries}p{3.2cm} >{\raggedright\arraybackslash}X >{\raggedright\arraybackslash}X}
\toprule
Nivel & Qué hace & Límite \\
\midrule
Nación (Congreso) & Sanciona la ley de presupuestos mínimos ambientales: una base común. & Fija el piso de protección. \\
Provincias & Dictan normas en sus jurisdicciones para complementar y aumentar los requisitos de protección. & No pueden estar por debajo de los presupuestos mínimos. \\
\bottomrule
\end{tabularx}
\end{table}

\section{Principios de la ley de presupuestos mínimos}

La ley de presupuestos mínimos ambientales contiene principios, que se desarrollan a continuación.

\subsection{Principio de prevención}

La ley tiene como objeto primordial anticiparse a la concreción de un daño ambiental y prevenirlo. Como el daño ambiental es un daño complejo de reconstruir o remediar y afecta el interés general, todas las acciones deben estar destinadas a prevenir los efectos negativos que una actividad puede tener sobre el ambiente.

\subsection{Principio precautorio}

Muy ligado al anterior, el principio precautorio exige que, ante un peligro de daño grave, prime una visión prudencial que tienda a impedir la degradación del ambiente y a tomar medidas. Está destinado a las autoridades administrativas y también a los jueces. Cuando los jueces ven, en el marco de un amparo ambiental o de una pretensión ambiental, que puede haber un peligro, debe primar la idea de precaución. Frente a un peligro cierto, no ante una mera conjetura, es decir, cuando hay una fuerte posibilidad de concreción del daño, los jueces deberían ser propensos a dictar medidas cautelares para evitar la irreparabilidad del daño.

\subsection{Principio de responsabilidad}

La responsabilidad incumbe al responsable del daño o del peligro ambiental, y debe ser enteramente satisfactoria para la finalidad de la ley: el daño ambiental debe remediarse hasta el máximo posible. Quien genera el daño debe suspender de forma inminente la acción que lo produce y tender a la \textbf{reparación integral} del daño ambiental.

\begin{table}[H]
\centering
\small
\caption{Principios ambientales}
\begin{tabularx}{\textwidth}{>{\raggedright\arraybackslash\bfseries}p{2.6cm} >{\raggedright\arraybackslash}X >{\raggedright\arraybackslash}X}
\toprule
Principio & Contenido & Destinatarios / aplicación \\
\midrule
Prevención & Anticiparse a la concreción del daño, que es complejo de reconstruir y afecta el interés general. & Todas las acciones deben dirigirse a prevenir los efectos negativos de una actividad. \\
Precautorio & Ante un peligro de daño grave y cierto (no una mera conjetura), prima la prudencia para impedir la degradación. & Autoridades administrativas y jueces: propensión a dictar medidas cautelares en amparos ambientales. \\
Responsabilidad & Quien genera el daño debe suspender la acción dañosa y tender a la reparación integral. & El responsable del daño o del peligro ambiental. \\
\bottomrule
\end{tabularx}
\end{table}

\subsection{Equidad intergeneracional}

\begin{definition}{Equidad intergeneracional}
Principio referido, sobre todo, a la explotación de los recursos naturales: en los recursos agotables y no renovables debe primar una política de gobierno tendiente a un uso racional y a la protección de las generaciones futuras. La explotación debe hacerse de forma ambientalmente responsable, de modo de no agotar irracionalmente los recursos ni destruirlos cuando son renovables.
\end{definition}

El ambiente no es dado por las generaciones previas, sino que, como se dice comúnmente, «nos es prestado por las futuras generaciones». De allí surge el deber de uso racional de los recursos, para que no se agoten y las nuevas generaciones puedan acceder también a ellos. Es un concepto general, cuyo cumplimiento debe luego verse en cada caso.

\begin{example}{Vaca Muerta}
Argentina ha iniciado, sobre todo en la gestión actual pero incluso desde mucho antes, la explotación de hidrocarburos no convencionales, como Vaca Muerta. La explotación debe ser racional y causar el menor daño posible al ambiente. Aun cuando son zonas alejadas de los centros urbanos, la degradación del subsuelo es inevitable y debe realizarse de forma racional.
\end{example}

El principio aplica tanto a los recursos no reproducibles, como los minerales, como a los que tienen capacidad de regeneración, como los suelos y las cuencas hídricas, que pueden degradarse si se explotan de manera constante e irracional. Se mencionaron como ejemplos:
\begin{itemize}
    \item la pesca en el mar argentino, donde hay acciones tendientes a establecer cotos de pesca;
    \item la agricultura, que requiere cuidados como el abono de la tierra y evitar el abuso de determinados insumos, como el glifosato o cualquiera de los paquetes tecnológicos utilizados en la siembra.
\end{itemize}
El entramado de normas es muy extenso. El concepto de equidad intergeneracional tiende a la explotación racional de los recursos, de forma tal que permita, para las futuras generaciones, una reserva y una fuente de aprovisionamiento en los años venideros.

\section{Información, evaluación de impacto y leyes sectoriales}

\subsection{Acceso a la información}

En materia ambiental es muy importante el acceso a la información. Debe existir la total posibilidad de acceder a información clara, seria y fehaciente, porque la materia ambiental involucra a todos como cohabitantes del mismo suelo; por ello existe siempre un deber de información claro.

\subsection{Evaluación de impacto ambiental}

La evaluación de impacto ambiental es una técnica preventiva que se ejecuta cuando hay obras relevantes con impacto en el ambiente. Todos los gobiernos, nacionales y subnacionales, deben encararla con responsabilidad, de modo que cada construcción o intervención relevante esté precedida de un análisis cierto, serio y profesional de los posibles efectos de la obra o explotación, y que exista un \textbf{plan de mitigación y remediación}.

Frente a cada acción relevante con efecto sobre el ambiente debe haber, como preludio, un análisis que se da en la fase administrativa: quien quiera hacer la obra debe analizar los impactos y contar con un plan tendiente a mitigar los efectos dañinos sobre el ambiente y la remediación.

\subsection{Ley general y leyes sectoriales}

La ley de presupuestos mínimos prevé muchas cuestiones en su texto originario, que luego se retroalimentan con otras leyes específicas: leyes contra incendios forestales, de preservación de bosques nativos, del uso del suelo y sobre el aprovechamiento de los recursos hídricos. Existe una suerte de ley general que se compagina con leyes sectoriales más específicas, que coparticipan de los mismos principios clave.

\section{Remediación, cooperación y legitimación ambiental}

\subsection{Remediación y seguros ambientales}

La remediación está a cargo de quien genera un eventual daño ambiental y debe ser integral. La ley prevé también el desarrollo de seguros ambientales, que permitan costear los eventuales daños y poner en práctica las labores de remediación.

\subsection{Cooperación y solidaridad: el daño interjurisdiccional}

Existen conceptos de cooperación y solidaridad entre la Nación y las provincias. Cooperación significa que todos los organismos ambientales deben intercambiar información y trabajar en conjunto, porque el ambiente es, por su propia naturaleza, algo inasible e interjurisdiccional: lo que se emite desde un establecimiento fabril a la atmósfera, o lo que se vierte en un río, va más allá del límite artificial entre la Capital Federal y la provincia de Buenos Aires, y un río puede bañar varias jurisdicciones, como la cuenca del Riachuelo.

\begin{example}{Cuenca del Riachuelo}
En un fallo de la Corte intervinieron la Nación, la Ciudad de Buenos Aires, la Provincia y numerosos municipios. Por su propia naturaleza, el daño ambiental tiende a ser interjurisdiccional, por lo que se requiere cooperación en todos los aspectos de la protección ambiental entre los organismos ambientales de la Nación, las provincias e incluso los municipios.
\end{example}

\subsection{Máxima amplitud en la legitimación activa}

Frente a la interjurisdiccionalidad del daño, y para propender a una fuerte labor de control y seguimiento, la Constitución y la jurisprudencia de la Corte han destacado que, para los reclamos judiciales ambientales, debe darse el \textbf{máximo nivel de apertura en la legitimación activa}. Pueden reclamar:
\begin{itemize}
    \item el defensor del pueblo, de la jurisdicción que sea o de la Nación;
    \item las asociaciones de defensa ambiental;
    \item los grupos de damnificados.
\end{itemize}

\section{Cuadro Cornell: medio ambiente}

\begin{longtable}{
    >{\raggedright\arraybackslash}p{\cornellizq}
    >{\raggedright\arraybackslash}p{\cornellder}
}
\toprule
\textbf{Preguntas / palabras clave} & \textbf{Notas / respuestas} \\
\midrule
\endhead

\textbf{¿Quién regula el ambiente según la Constitución?} &
La Nación dicta los presupuestos mínimos, que constituyen un piso que no se puede perforar, y las provincias los complementan elevando la protección dentro de su jurisdicción. Rige la complementariedad. \\[0.8em]

\textbf{¿Qué establecen los principios de prevención, precaución y responsabilidad?} &
\textbf{Prevención:} anticiparse al daño, que es complejo de reconstruir y afecta el interés general. \textbf{Precautorio:} ante un peligro cierto de daño grave, tomar medidas; los jueces pueden dictar cautelares. \textbf{Responsabilidad:} quien daña debe cesar la conducta y reparar integralmente. \\[0.8em]

\textbf{¿Qué significa equidad intergeneracional?} &
Uso racional de los recursos naturales (sobre todo los no renovables) para no agotarlos y proteger a las generaciones futuras: el ambiente «nos es prestado por las futuras generaciones». Ejemplos: Vaca Muerta, pesca, agricultura. \\[0.8em]

\textbf{¿Qué rol cumplen la información y la evaluación de impacto ambiental?} &
Debe garantizarse el acceso a información clara, seria y fehaciente. Toda obra o explotación relevante necesita un análisis serio previo de sus efectos y un plan de mitigación y remediación. La ley general se complementa con leyes sectoriales. \\[0.8em]

\textbf{¿Quién repara el daño ambiental?} &
Quien lo genera; la reparación debe ser integral y la ley prevé seguros ambientales para costear daños y remediación. \\[0.8em]

\textbf{¿Por qué es necesaria la cooperación y quién puede reclamar judicialmente?} &
Porque el daño es interjurisdiccional (caso del Riachuelo), por lo que hay cooperación entre Nación, provincias y municipios. La legitimación activa es amplia: defensor del pueblo, asociaciones y damnificados. \\

\midrule
\multicolumn{2}{p{\dimexpr\cornellizq+\cornellder+2\tabcolsep\relax}}{%
\textbf{Resumen:} Todos tienen derecho a un ambiente sano y el daño ambiental genera deber de reparación. La Nación fija un piso de presupuestos mínimos que las provincias pueden mejorar. Rigen los principios de prevención, precaución, responsabilidad y equidad intergeneracional, con acceso a la información, evaluación de impacto con plan de mitigación y remediación, remediación integral a cargo de quien daña, seguros ambientales, cooperación entre jurisdicciones y legitimación activa de máxima amplitud. } \\
\bottomrule
\end{longtable}

% ============================================================
% CAPÍTULO 3
% ============================================================
\chapter{Derecho de propiedad y expropiación}

\section{Artículo 17: inviolabilidad de la propiedad}

La consagración más clara del derecho de propiedad a nivel constitucional se encuentra en el artículo 17, que establece que la propiedad es inviolable.

\subsection{Qué significa «inviolable» y qué se entiende por propiedad}

La inviolabilidad de la propiedad opera tanto frente a \textbf{terceros} que puedan turbar el uso y goce de la propiedad como frente al \textbf{propio Estado}, nacional o provincial.

\begin{definition}{Propiedad}
Conjunto de bienes, tangibles o intangibles, que son bienes jurídicos tutelados y que se diferencian del propio cuerpo del sujeto que los tiene. Abarca bienes inmuebles, dinero, valores y derechos jurídicos, como un título de propiedad intelectual. Una pieza musical compuesta, una idea registrada y protegida por la propiedad intelectual o el uso de una marca que puede explotarse o venderse forman parte del patrimonio.
\end{definition}

La Constitución consagra, entonces, un concepto de inviolabilidad destinado a los terceros y también a la acción de gobierno. La ley de expropiación es la única herramienta jurídica que permite un desapoderamiento de bienes o de la propiedad legítima.

\begin{important}
La inviolabilidad de la propiedad opera frente a terceros y frente al Estado (nacional o provincial). La única herramienta jurídica que permite desapoderar legítimamente a una persona de sus bienes es la ley de expropiación.
\end{important}

\section{Requisitos y alcance de la expropiación}

La facultad expropiatoria no constituye un cheque en blanco para la administración. La ley de expropiación contiene numerosas regulaciones.

\subsection{Indemnización previa y declaración de utilidad pública}

\begin{itemize}
    \item \textbf{Indemnización previa:} el pago debe ser previo al desapoderamiento del bien. La Constitución también lo dice.
    \item \textbf{Intervención del Congreso:} el Poder Ejecutivo no puede expropiar un bien sin una declaración previa de utilidad pública por parte del Poder Legislativo. Es un punto clave del sistema republicano de división de poderes: el Ejecutivo no puede avanzar si no parte de una ley de declaración de utilidad pública, y puede avanzar en el cumplimiento de la expropiación, pero no en el desapoderamiento, porque el pago tiene que ser previo.
    \item \textbf{Determinación de los bienes:} los bienes tienen que estar determinados por la ley. No puede haber una ley indefinida que diga que el Estado puede expropiar todos los terrenos que considere de utilidad pública. Los bienes deben estar determinados taxativamente, o al menos con suficiente claridad, para que el Ejecutivo se ajuste a una hoja de ruta concreta.
\end{itemize}

\subsection{¿Qué debe indemnizar el Estado?}

El Estado debe indemnizar el \textbf{valor objetivo del bien}, con prescindencia del lucro cesante, de determinado daño moral y de valores sentimentales. Es lo que el bien cuesta objetivamente en el mercado. Se indemnizan también los perjuicios directos vinculados con el desapoderamiento. Por lo tanto, no es una indemnización plena.

\subsection{El proyecto de ley sobre protección de la propiedad privada}

Se mencionó un proyecto de ley de la gestión actual sobre protección de la propiedad privada, que pretende ampliar el rango de indemnización del bien, para que el Estado también pague, hasta un punto. Según lo expuesto, el proyecto prevé indemnizar también el lucro cesante, pero hasta un determinado límite, expresado no como un monto cuantitativo sino como un porcentaje del daño emergente. Al momento de la clase estaba en debate en el Congreso una norma que ampliaría la indemnización que el Estado debería pagar por expropiación.

% TODO: contenido incompleto o ambiguo en las notas originales (el contenido del proyecto de ley se expone de manera aproximada: «creo que el proyecto de ley habla de...»)

Sin una ley previa que declare la utilidad pública, el Estado no puede desapoderar a los particulares de sus bienes.

\begin{table}[H]
\centering
\small
\caption{Requisitos y alcance de la expropiación}
\begin{tabularx}{\textwidth}{>{\raggedright\arraybackslash\bfseries}p{3.4cm} >{\raggedright\arraybackslash}X}
\toprule
Requisito & Contenido \\
\midrule
Quién declara la utilidad pública & El Congreso, mediante ley. El Poder Ejecutivo no puede expropiar sin esa declaración previa (división de poderes). \\
Determinación de los bienes & La ley debe determinarlos taxativamente o con suficiente claridad; no hay «cheque en blanco» para la administración. \\
Momento de la indemnización & Previa al desapoderamiento: sin pago previo no hay desapoderamiento. \\
Qué se indemniza & El valor objetivo del bien (lo que cuesta en el mercado) y los perjuicios directos vinculados al desapoderamiento. \\
Qué no se indemniza & Lucro cesante, daño moral y valores sentimentales: no es una indemnización plena. \\
\bottomrule
\end{tabularx}
\end{table}

\section{Procedimiento y expropiación inversa}

\subsection{El avenimiento y el juicio de expropiación}

El procedimiento de la ley de expropiación comienza con una etapa de acuerdo, denominada \textbf{avenimiento}. Si el particular acepta la oferta que hace el Estado por el bien, puede ser retribuido hasta en un 10\,\% del valor de la tasación que haga el Tribunal de Tasaciones, como incentivo para que se avenga con la administración y entregue su bien a cambio de un precio convenido.
% TODO: contenido incompleto o ambiguo en las notas originales (no se precisa si el 10 % es un adicional sobre la tasación)

Si frente al expropiante, Estado nacional o provincial, la parte no está de acuerdo con el precio ofrecido (por ejemplo, porque considera que el bien está mal tasado o fuera de mercado), puede oponerse y se genera el \textbf{juicio de expropiación}. En él, las dos partes actúan como contrapartes y el juez resuelve el monto de la indemnización, considerada objetiva y justa.

\subsection{Expropiación ordinaria y expropiación inversa}

El particular tiene la garantía de que, si considera que lo despojan injustamente de su bien porque no le pagan lo que corresponde, puede reclamarlo, como actor o como demandado, porque la ley de expropiaciones prevé la expropiación inversa.

\textbf{Expropiación ordinaria.} Es la típica: el Estado, luego de que el Congreso dicta la ley de utilidad pública, realiza una tasación e inicia, frente a la oposición del particular que no se aviene, un juicio de expropiación para desapoderarlo del bien.

\begin{definition}{Expropiación inversa}
Supuesto en el cual el propio particular acude al juez y solicita que se expropie su bien.
\end{definition}

La expropiación inversa procede en dos supuestos:
\begin{enumerate}
    \item \textbf{Inactividad estatal:} la ley de expropiación se dictó pero no se puso en ejecución. El particular pide que se ejecute la ley; por ejemplo, porque estando la ley sancionada nadie querrá comprarle su departamento y desea venderlo, por lo que intenta que el Estado se lo expropie.
    \item \textbf{Imposibilidad de uso y goce razonable:} sobre el bien de su propiedad existe algún tipo de regulación o actividad estatal que genera que el particular no pueda usar y gozar de esa propiedad razonablemente. El bien resulta inviable para ser usufructuado y debe invocarse que el Estado expropie.
\end{enumerate}

\begin{example}{Supuestos de expropiación inversa}
\begin{itemize}
    \item Un negocio al que el Estado traza una autopista o modifica una ruta que pasa por delante del establecimiento, de modo que no se puede entrar ni salir. Aunque el inmueble no fue expropiado, la autopista tan cercana impide utilizarlo.
    \item Una represa que inunda un camino que daba acceso a un campo, de modo que ya no se puede acceder a él.
\end{itemize}
\end{example}

\begin{table}[H]
\centering
\small
\caption{Expropiación ordinaria y expropiación inversa}
\begin{tabularx}{\textwidth}{>{\raggedright\arraybackslash\bfseries}p{2.4cm} >{\raggedright\arraybackslash}X >{\raggedright\arraybackslash}X}
\toprule
Aspecto & Expropiación ordinaria & Expropiación inversa \\
\midrule
Quién la promueve & El Estado (nacional o provincial). & El propio particular, que acude al juez. \\
Punto de partida & Ley de utilidad pública, tasación y oferta; si el particular no se aviene, juicio de expropiación. & Dos supuestos: (1) la ley de expropiación se dictó pero el Estado no la ejecuta (inactividad estatal); (2) una regulación o actividad estatal impide usar y gozar razonablemente del bien. \\
Ejemplos & Inmueble tasado por el Tribunal de Tasaciones. & Departamento incluido en una ley no ejecutada; negocio al que una autopista impide entrar y salir; campo cuyo camino de acceso inunda una represa. \\
Resolución & El juez fija el monto de la indemnización objetiva y justa. & El propietario pide: «Estado, exprópieme este bien». \\
\bottomrule
\end{tabularx}
\end{table}

\begin{note}
Frente a la inquietud de si, en los supuestos de imposibilidad de uso y goce por una actividad estatal, puede reclamarse daños y perjuicios por responsabilidad estatal por actividad lícita en lugar de pedir la expropiación, se explicó que, si no se puede usar y gozar de la propiedad por algún tipo de actividad estatal, no puede pedirse daños y perjuicios: debe pedirse la expropiación del bien.
\end{note}

\section{Cuadro Cornell: derecho de propiedad y expropiación}

\begin{longtable}{
    >{\raggedright\arraybackslash}p{\cornellizq}
    >{\raggedright\arraybackslash}p{\cornellder}
}
\toprule
\textbf{Preguntas / palabras clave} & \textbf{Notas / respuestas} \\
\midrule
\endhead

\textbf{¿Dónde se consagra el derecho de propiedad y qué abarca?} &
En el artículo 17 de la Constitución: la propiedad es inviolable frente a terceros y frente al Estado. Abarca bienes tangibles e intangibles (inmuebles, dinero, valores, derechos, propiedad intelectual, marcas). \\[0.8em]

\textbf{¿Cuál es la única herramienta que permite desapoderar legítimamente de un bien?} &
La ley de expropiación. \\[0.8em]

\textbf{¿Qué exige la expropiación?} &
Ley del Congreso que declare la utilidad pública y determine los bienes (el Ejecutivo no puede expropiar sin ella), e indemnización previa al desapoderamiento. \\[0.8em]

\textbf{¿Qué se indemniza y qué no?} &
Se indemniza el valor objetivo del bien y los perjuicios directos; no se indemniza lucro cesante, daño moral ni valores sentimentales. Un proyecto en debate busca ampliar el alcance. \\[0.8em]

\textbf{¿Cómo es el procedimiento expropiatorio?} &
Primero el avenimiento (retribución de hasta un 10\,\% del valor de la tasación del Tribunal de Tasaciones); si no hay acuerdo, juicio de expropiación, donde el juez fija el monto de la indemnización objetiva y justa. \\[0.8em]

\textbf{¿Qué es la expropiación inversa y cuándo procede?} &
El propietario pide ser expropiado cuando hay inactividad estatal (ley dictada pero no ejecutada) o cuando una regulación o actividad estatal (autopista, represa) le impide usar y gozar razonablemente su bien. \\

\midrule
\multicolumn{2}{p{\dimexpr\cornellizq+\cornellder+2\tabcolsep\relax}}{%
\textbf{Resumen:} El artículo 17 consagra la inviolabilidad de la propiedad, frente a terceros y frente al Estado. Solo puede limitarse mediante la expropiación, con ley previa del Congreso que declare la utilidad pública y determine los bienes, e indemnización previa del valor objetivo del bien y los perjuicios directos. El procedimiento tiene una etapa de avenimiento y, si no hay acuerdo, un juicio de expropiación. El propietario puede activar la expropiación inversa por inactividad estatal o cuando no puede usar y gozar razonablemente su bien. } \\
\bottomrule
\end{longtable}

% ============================================================
% CAPÍTULO 4
% ============================================================
\chapter{Responsabilidad del Estado y contratos}

\section{Responsabilidad estatal por actividad legítima}

La expropiación es una suerte de responsabilidad por actividad legítima. El Estado se condujo con apego a la legalidad: el Legislativo dictó la ley de utilidad pública y el Ejecutivo la ejecutó. Desapoderó del bien, pero antes pagó lo que correspondía. Es una actividad legítima, porque está prevista en la ley de expropiaciones.

\subsection{Cambio normativo y derecho adquirido}

Otro supuesto de actividad legítima que responsabiliza al Estado se da cuando se modifica una ley o norma y existía un \textbf{derecho adquirido} al amparo de la ley que se deroga. El Estado puede cambiar la ley, porque nadie tiene derecho a la eterna duración de las normas: nadie tiene derecho al mantenimiento de las leyes y los reglamentos. Pero quien tenía un derecho amparado en la ley que luego se cambia sí tiene derecho a una indemnización.

\begin{example}{Habilitación para abrir un negocio}
Una ley autorizaba a abrir un negocio y luego la administración, por razones de utilidad pública, entiende que no se puede, y retira esa habilitación. El titular no puede resistir el cambio legislativo, pero debe ser indemnizado por lo que perdió, que era un derecho concreto a comercializar un determinado bien.
\end{example}

\begin{example}{Estaciones de servicio bajo edificios}
Hubo un caso famoso: las estaciones de servicio no podían estar debajo de edificios, es decir, no podía haber edificaciones por encima de la estación, por una cuestión de seguridad. La norma salió en los años 70, pero mucho tiempo antes había estaciones de servicio habilitadas bajo edificios. Quien tenía una estación con pisos edificados arriba debió cerrarla. El cierre se funda en una regulación por el interés común de preservar la salud de la gente y evitar incendios y explosiones. Pero el titular tenía una estación funcionando con todas las de la ley; producto del cambio legal quedó imposibilitado de mantener su negocio y debe ser indemnizado para mitigar los daños, porque tenía un derecho adquirido bajo la ley que luego se modificó.
\end{example}

\begin{example}{Planta de baterías de litio (ejemplo hipotético)}
Supóngase que se descubre que las baterías de litio son gravosas para la salud. Una empresa tenía una concesión para el litio y se le niega mantener su planta generadora de baterías por razones ambientales. La prohibición posterior es legítima porque está fundada en el derecho ambiental, pero la empresa tenía un derecho adquirido a esa comercialización y a que su establecimiento industrial estuviera habilitado, y debe ser resarcida.
\end{example}

\begin{important}
Nadie tiene derecho al mantenimiento de las leyes y reglamentos, pero quien había consolidado un derecho adquirido bajo la norma derogada sí tiene derecho a ser indemnizado: el Estado puede cambiar la ley, pero debe resarcir el derecho concreto que se pierde.
\end{important}

\section{Responsabilidad estatal por actividad ilegítima}

\begin{definition}{Falta de servicio}
Supuesto en el cual el Estado tenía que cumplir con un cometido, con un servicio o con una determinada competencia, y no lo hizo. Es la base de la responsabilidad estatal por actividad ilegítima.
\end{definition}

En este caso el Estado no se conduce con apego a la ley. Debe indemnizar \textbf{integralmente}: a diferencia de la expropiación o de la responsabilidad por actividad legítima, en las que la indemnización es reducida (en la expropiación, solo daño emergente), el Estado cubre todos los daños, incluyendo pérdida de chance, lucro cesante, daño moral, etcétera.

\begin{example}{Guardia de un hospital público}
Una persona llega a la guardia de un hospital público para ser atendida y se demuestra una negligencia en el servicio: se tardó mucho en el traslado y luego en el ingreso a la sala de operaciones, y el paciente fallece o queda con una incapacidad. El Estado, que tenía el deber de cuidar la salud de esa persona y la preservación de ese bien de vida, puede ser sancionado con el pago de una indemnización, que no es reducida como en la actividad lícita, sino integral.
\end{example}

En el caso de la actividad legítima, el Estado no tuvo un obrar ilegítimo ni actuó ilícitamente; actuó en el marco de una competencia, como la expropiación o el cambio legislativo, pero dañó de forma específica un derecho adquirido o la propiedad.

\begin{table}[H]
\centering
\small
\caption{Responsabilidad del Estado por actividad legítima e ilegítima}
\begin{tabularx}{\textwidth}{>{\raggedright\arraybackslash\bfseries}p{2.4cm} >{\raggedright\arraybackslash}X >{\raggedright\arraybackslash}X}
\toprule
Aspecto & Actividad legítima (lícita) & Actividad ilegítima (ilícita) \\
\midrule
Conducta del Estado & Actúa con apego a la legalidad, en el marco de una competencia. & Falta de servicio: debía cumplir un cometido o competencia y no lo hizo. \\
Ejemplos & Expropiación; cambio normativo que afecta un derecho adquirido (estaciones de servicio; planta de litio). & Negligencia en la guardia de un hospital público (demora en el traslado y en el ingreso a quirófano). \\
Qué se daña & Un derecho adquirido o la propiedad. & Bienes y derechos de la persona, incluida la vida y la salud. \\
Indemnización & Reducida: mitiga el daño del derecho concreto perdido (en la expropiación, daño emergente). & Integral: todos los daños, incluyendo pérdida de chance, lucro cesante y daño moral. \\
\bottomrule
\end{tabularx}
\end{table}

\section{Propiedad y contratos}

La Corte ha reconocido el derecho de propiedad, como traslación del artículo 17, a los contratos. Los contratos son ley para las partes y el Estado no debería entrometerse ni modificarlos unilateralmente, ya sea un contrato entre terceros o un contrato que el particular tenía con el Estado. Los contratos son fuente originadora de derechos y obligaciones.

\begin{example}{Provisión de resmas a la Universidad de Buenos Aires}
Un proveedor vendía resmas a la Universidad de Buenos Aires para sus impresiones. Si el Estado modifica el contrato, el incumplimiento contractual afecta el derecho de propiedad del proveedor, que puede reclamar al Estado nacional, o a quien corresponda, el daño producido.
\end{example}

En aras del bien común, el Estado no puede romper el contrato ni tornarlo más oneroso para la otra parte. Eso debería ser motivo de una compensación o indemnización, para compatibilizar el derecho de propiedad con el interés público.

\section{Cuadro Cornell: responsabilidad del Estado y contratos}

\begin{longtable}{
    >{\raggedright\arraybackslash}p{\cornellizq}
    >{\raggedright\arraybackslash}p{\cornellder}
}
\toprule
\textbf{Preguntas / palabras clave} & \textbf{Notas / respuestas} \\
\midrule
\endhead

\textbf{¿Cuándo responde el Estado por actividad legítima?} &
Cuando actúa conforme a derecho (expropiación o cambio normativo) pero afecta un derecho adquirido. Nadie tiene derecho al mantenimiento de las leyes, pero sí a ser indemnizado por el derecho concreto perdido (estaciones de servicio, planta de litio). La indemnización es reducida. \\[0.8em]

\textbf{¿Qué es la falta de servicio y cuánto debe pagar el Estado?} &
El Estado tenía que cumplir un cometido y no lo hizo (ejemplo del hospital público). La indemnización es integral: incluye pérdida de chance, lucro cesante y daño moral. \\[0.8em]

\textbf{¿Qué diferencia hay entre la indemnización por actividad legítima e ilegítima?} &
En la legítima es reducida (en la expropiación, daño emergente); en la ilegítima es integral y cubre todos los daños. \\[0.8em]

\textbf{¿Cómo se vincula la propiedad con los contratos?} &
La Corte extendió el artículo 17 a los contratos: son ley para las partes y fuente de derechos y obligaciones. El Estado no puede romperlos ni volverlos más onerosos unilateralmente; si lo hace, corresponde compensación para compatibilizar propiedad e interés público. \\

\midrule
\multicolumn{2}{p{\dimexpr\cornellizq+\cornellder+2\tabcolsep\relax}}{%
\textbf{Resumen:} El Estado responde por su actividad legítima (expropiación, cambio normativo que afecta derechos adquiridos) con una indemnización reducida, y por su actividad ilegítima (falta de servicio) con una indemnización integral. Nadie tiene derecho al mantenimiento de las leyes, pero sí a ser resarcido por un derecho adquirido que se pierde. Los contratos, protegidos por el derecho de propiedad, no pueden ser alterados unilateralmente por el Estado sin compensación. } \\
\bottomrule
\end{longtable}

% ============================================================
% CAPÍTULO 5
% ============================================================
\chapter{Legitimación activa: el caso Halabi}

\section{Caso Halabi: tres categorías de legitimación}

\begin{definition}{Legitimación activa}
Aptitud para ser parte en un juicio como parte actora, es decir, para poder demandar judicialmente una pretensión.
\end{definition}

En el precedente Halabi, que tiene ya bastante recorrido, la Corte Suprema señaló tres categorías o presupuestos de la legitimación activa: los derechos subjetivos, los derechos de incidencia colectiva sobre bienes colectivos y los derechos de incidencia colectiva sobre intereses individuales homogéneos.

\subsection{Primera categoría: derechos de incidencia colectiva sobre bienes colectivos}

Es el concepto más general. Por ejemplo, en el tema ambiental todos tienen derecho a un ambiente sano y un eventual daño ambiental puede perjudicar a todos los que están en una misma situación geográfica. Son muchos los sujetos que pueden iniciar acciones judiciales y pedir la protección de ese derecho, que es afectado de forma colectiva.

\subsection{Segunda categoría: intereses individuales homogéneos}

Es una categoría intermedia, de derechos colectivos. No es colectivo porque todos sean copartícipes del derecho, sino porque se trata de un daño originado sobre una pluralidad de sujetos damnificados.

\begin{example}{25.000 camionetas con falla de frenos}
En una planta de montaje de automóviles, durante tres meses, por una mala funcionalidad o un error, se produjeron (con números supuestos muy altos) 25.000 camionetas con un problema en los frenos, que deben volver al taller para su revisión y reparación. Hay una pluralidad de sujetos, todos los que compraron esas camionetas, que pueden ser resarcidos. Es un número concreto y determinado, en función de un mismo hecho dañoso que se despliega sobre los bienes de todos ellos, con una determinación específica del bien o de la situación a indemnizar. La homogeneidad está dada porque todos los casos se originaron en una misma circunstancia: el error en la planta de montaje.
\end{example}

La legitimación activa corresponde a todos los damnificados, pero también puede llevar adelante un eventual reclamo una asociación de defensa del consumidor o el defensor del pueblo. El defensor del pueblo participa en la medida en que no haya otro sector social amenazado o demandado, porque no puede arbitrar entre los intereses de una sociedad.

\begin{example}{El defensor del pueblo y el corralito}
Con motivo del corralito, en materia de préstamos, el defensor del pueblo quería declarar la nulidad de la pesificación de los contratos de préstamo. Pero esto tenía vencedores y vencidos: a algunos les convenía y a otros no. Por ello, la legitimación de las asociaciones de defensa del consumidor debe ser siempre en defensa de esa parte débil, el consumidor.
\end{example}

\subsection{Tercera categoría: derechos subjetivos}

Los derechos subjetivos son los que cada persona detenta con individualidad y excluyendo generalmente a los demás. Quien tiene un derecho de propiedad sobre su vivienda, con registro de la propiedad, escritura traslativa y posesión, puede excluir de ese derecho a todos los que intenten ocupar la vivienda. Es un derecho típicamente individual porque excluye al resto. En cambio, en las otras dos figuras hay una categoría o grupo de personas que pueden tener el mismo derecho.

\begin{table}[H]
\centering
\small
\caption{Las tres categorías del caso Halabi}
\begin{tabularx}{\textwidth}{>{\raggedright\arraybackslash\bfseries}p{2.7cm} >{\raggedright\arraybackslash}X >{\raggedright\arraybackslash}p{2.9cm} >{\raggedright\arraybackslash}p{2.9cm}}
\toprule
Categoría & Características & Ejemplo & Quiénes pueden reclamar \\
\midrule
Derechos subjetivos & Individuales y excluyentes: cada uno los detenta con individualidad y excluye al otro. & Derecho de propiedad sobre la vivienda (escritura, registro, posesión). & El titular del derecho. \\
Incidencia colectiva sobre bienes colectivos & Concepto más general; el derecho pertenece a todos y el daño afecta a todos los que están en la misma situación. & Daño ambiental: todos tenemos derecho a un ambiente sano. & Defensor del pueblo, asociaciones y cualquier damnificado en representación de todos. \\
Incidencia colectiva sobre intereses individuales homogéneos & Pluralidad determinada de damnificados por un mismo hecho dañoso; la homogeneidad la da una misma causa. & 25.000 camionetas con falla de frenos fabricadas en tres meses por un error en la planta. & Los damnificados, asociaciones de defensa del consumidor y defensor del pueblo (si no hay otros sectores sociales amenazados). \\
\bottomrule
\end{tabularx}
\end{table}

\begin{important}
El defensor del pueblo no puede arbitrar entre intereses de la sociedad: cuando hay «vencedores y vencidos», la legitimación de las asociaciones de consumidores debe ser siempre en defensa de la parte débil, el consumidor.
\end{important}

\section{Acciones procesales según el tipo de derecho}

\textbf{Derechos de incidencia colectiva sobre bienes colectivos: amparo ambiental.} Puede promoverlo el defensor del pueblo, una asociación de defensa del ambiente o cualquier damnificado en representación de todos. La pretensión típica es hacer cesar la conducta que genera el daño ambiental.

\textbf{Intereses individuales homogéneos: acción de clase.} Puede promoverla el damnificado en representación del resto de los damnificados, o una asociación de defensa del consumidor en representación de todo el grupo de consumidores damnificado.

\begin{example}{Acción de clase por las camionetas}
Quien compra la camioneta inicia juicio contra la terminal automotriz y puede invocar que todos los que compraron esa camioneta durante el período indicado son también integrantes de esa clase.
\end{example}

\begin{table}[H]
\centering
\small
\caption{Acciones y pretensiones según el tipo de derecho}
\begin{tabularx}{\textwidth}{>{\raggedright\arraybackslash\bfseries}p{3.1cm} >{\raggedright\arraybackslash}p{2.2cm} >{\raggedright\arraybackslash}X >{\raggedright\arraybackslash}X}
\toprule
Tipo de derecho & Acción típica & Quién puede promoverla & Pretensión típica \\
\midrule
Incidencia colectiva sobre bienes colectivos & Amparo ambiental & Defensor del pueblo; asociación de defensa del ambiente; cualquier damnificado en representación de todos. & Hacer cesar la conducta que genera el daño ambiental. \\
Incidencia colectiva sobre intereses individuales homogéneos & Acción de clase & Un damnificado en representación del resto; asociación de defensa del consumidor. & Reparación a favor de todos los integrantes de la clase (por ejemplo, compradores de la camioneta en el período indicado). \\
\bottomrule
\end{tabularx}
\end{table}

\section{Cuadro Cornell: legitimación activa}

\begin{longtable}{
    >{\raggedright\arraybackslash}p{\cornellizq}
    >{\raggedright\arraybackslash}p{\cornellder}
}
\toprule
\textbf{Preguntas / palabras clave} & \textbf{Notas / respuestas} \\
\midrule
\endhead

\textbf{¿Qué es la legitimación activa?} &
La aptitud para ser parte en un juicio como parte actora, para poder demandar judicialmente una pretensión. \\[0.8em]

\textbf{¿Qué tres categorías de derechos fijó la Corte en Halabi?} &
\textbf{Derechos subjetivos} (individuales y excluyentes), \textbf{derechos de incidencia colectiva sobre bienes colectivos} (por ejemplo, el ambiente) y \textbf{derechos de incidencia colectiva sobre intereses individuales homogéneos} (por ejemplo, camionetas con falla de frenos). \\[0.8em]

\textbf{¿Qué caracteriza a los intereses individuales homogéneos?} &
Una pluralidad determinada de damnificados por un mismo hecho dañoso; la homogeneidad la da la causa común (el error en la planta de montaje). \\[0.8em]

\textbf{¿Por qué el defensor del pueblo no puede actuar cuando hay «vencedores y vencidos»?} &
Porque no puede arbitrar entre los intereses de una sociedad; la legitimación de las asociaciones de consumidores es siempre en defensa de la parte débil, el consumidor (ejemplo: pesificación de los préstamos en el corralito). \\[0.8em]

\textbf{¿Qué acción corresponde a cada categoría?} &
Bienes colectivos: amparo ambiental (defensor del pueblo, asociaciones o cualquier damnificado), con pretensión típica de cese del daño. Intereses homogéneos: acción de clase (un damnificado o una asociación de consumidores en representación del grupo). \\

\midrule
\multicolumn{2}{p{\dimexpr\cornellizq+\cornellder+2\tabcolsep\relax}}{%
\textbf{Resumen:} El caso Halabi ordena quién puede reclamar y cómo. Los derechos subjetivos son individuales y excluyentes; los derechos de incidencia colectiva sobre bienes colectivos pertenecen a todos (amparo ambiental, con pretensión de cese del daño); los derechos de incidencia colectiva sobre intereses individuales homogéneos agrupan a damnificados por una misma causa (acción de clase, promovida por un damnificado o por una asociación de consumidores). } \\
\bottomrule
\end{longtable}

% ============================================================
% SÍNTESIS FINAL
% ============================================================
\chapter*{Síntesis final}
\addcontentsline{toc}{chapter}{Síntesis final}
\markboth{Síntesis final}{Síntesis final}

Los cinco bloques comparten una misma idea de fondo: el poder, económico o estatal, tiene límites, y la Constitución ofrece herramientas para controlarlos y para reclamar cuando se los excede.

En \textbf{defensa de la competencia}, el Estado previene (control de concentraciones) y sanciona (multas) los abusos de quien domina un mercado, a través de una Autoridad independiente cuyas decisiones son apelables. En materia \textbf{ambiental}, la Nación fija un piso de presupuestos mínimos que las provincias pueden mejorar, y rigen los principios de prevención, precaución, responsabilidad y equidad intergeneracional, con evaluación de impacto, remediación integral y legitimación amplia. El \textbf{derecho de propiedad} del artículo 17 es inviolable frente a terceros y frente al Estado: solo puede ser limitado por expropiación, con ley previa de utilidad pública e indemnización previa del valor objetivo, y el propietario puede activar la expropiación inversa. Cuando el Estado daña, \textbf{responde}: de forma limitada si actuó legítimamente (cambio normativo sobre derechos adquiridos) e integralmente si incurrió en falta de servicio; y tampoco puede alterar unilateralmente los contratos. Finalmente, el caso \textbf{Halabi} ordena quién puede reclamar y cómo: derechos subjetivos, derechos de incidencia colectiva sobre bienes colectivos (amparo) y sobre intereses individuales homogéneos (acción de clase).

% ============================================================
% MATERIAL COMPLEMENTARIO
% ============================================================
\chapter{Material complementario: referencias normativas}

Las siguientes referencias figuran en el apunte original como complemento normativo. Se trata de textos agregados como referencia, a los que la clase alude por su contenido, sin citar necesariamente el número de artículo (salvo el artículo 17).

\begin{note}
\textbf{Constitución Nacional, art. 42 (2.º párrafo).} Las autoridades proveerán a la protección de esos derechos, a la educación para el consumo, a la defensa de la competencia contra toda forma de distorsión de los mercados, al control de los monopolios naturales y legales, al de la calidad y eficiencia de los servicios públicos, y a la constitución de asociaciones de consumidores y de usuarios.
\end{note}

\begin{note}
\textbf{Ley 27.442 --- Defensa de la Competencia.} Crea la Autoridad de Defensa de la Competencia, con dos competencias: una labor preventiva (control de concentraciones económicas) y una labor sancionatoria (multas por infracciones, por ejemplo el abuso de posición dominante). Sus decisiones son apelables ante el fuero civil y comercial federal; el órgano tiene cinco miembros designados con acuerdo del Senado. Síntesis de lo expuesto en clase; el articulado no fue leído.
\end{note}

\begin{note}
\textbf{Constitución Nacional, art. 41 (primeros tres párrafos).} Todos los habitantes gozan del derecho a un ambiente sano, equilibrado, apto para el desarrollo humano y para que las actividades productivas satisfagan las necesidades presentes sin comprometer las de las generaciones futuras; y tienen el deber de preservarlo. El daño ambiental generará prioritariamente la obligación de recomponer, según lo establezca la ley. Las autoridades proveerán a la protección de este derecho, a la utilización racional de los recursos naturales, a la preservación del patrimonio natural y cultural y de la diversidad biológica, y a la información y educación ambientales. Corresponde a la Nación dictar las normas que contengan los presupuestos mínimos de protección, y a las provincias, las necesarias para complementarlas, sin que aquéllas alteren las jurisdicciones locales.

La «Ley de Presupuestos Mínimos» a la que se alude es, en su versión general, la Ley General del Ambiente (n.º 25.675), según la referencia agregada en el apunte original.
\end{note}

\begin{note}
\textbf{Constitución Nacional, art. 17.} La propiedad es inviolable, y ningún habitante de la Nación puede ser privado de ella, sino en virtud de sentencia fundada en ley. La expropiación por causa de utilidad pública, debe ser calificada por ley y previamente indemnizada. Sólo el Congreso impone las contribuciones que se expresan en el art. 4º. Ningún servicio personal es exigible, sino en virtud de ley o de sentencia fundada en ley. Todo autor o inventor es propietario exclusivo de su obra, invención o descubrimiento, por el término que le acuerde la ley. La confiscación de bienes queda borrada para siempre del Código Penal argentino. Ningún cuerpo armado puede hacer requisiciones, ni exigir auxilios de ninguna especie.
\end{note}

\begin{note}
\textbf{Ley de Expropiación (Ley 21.499).} Referida en clase como «la ley de expropiación»: declaración de utilidad pública por ley, bienes determinados, indemnización previa del valor objetivo y perjuicios directos, etapa de avenimiento, juicio de expropiación y expropiación inversa. La identificación por su número es un agregado del apunte original; el articulado no fue leído.
\end{note}

\begin{note}
\textbf{Constitución Nacional, art. 43 (2.º párrafo) y caso Halabi.} Podrán interponer esta acción contra cualquier forma de discriminación y en lo relativo a los derechos que protegen al ambiente, a la competencia, al usuario y al consumidor, así como a los derechos de incidencia colectiva en general, el afectado, el defensor del pueblo y las asociaciones que propendan a esos fines, registradas conforme a la ley, la que determinará los requisitos y formas de su organización.

Corte Suprema de Justicia de la Nación, «Halabi, Ernesto c/ P.E.N. -- ley 25.873 -- dto. 1563/04 s/ amparo ley 16.986», 24 de febrero de 2009 (Fallos: 332:111).
\end{note}

\end{document}
